\documentclass[12pt]{article}
% Préambule commun à tous les documents extraits de « La Revue CAPES »
\usepackage[T1]{fontenc}
\usepackage[utf8]{inputenc}
\usepackage{lmodern}
\usepackage{amsmath,amssymb,amsthm,amsfonts,mathrsfs}
\usepackage{setspace}
\usepackage{listings}
\usepackage{pgfplots}
\pgfplotsset{compat=1.18}
\usepackage{longtable}
\usepackage{adjustbox}
\usepackage{lettrine}
\usepackage{tkz-tab}
\usepackage{changepage}
\usepackage{pdflscape}
\usepackage{graphicx}
\usepackage{tikz}
\usepackage{xcolor}
\usepackage[a4paper,top=2cm,bottom=2.5cm,left=2.5cm,right=2cm]{geometry}
\usepackage{fancyhdr}
\usepackage{draftwatermark}
\SetWatermarkText{La RevueCapes}
\SetWatermarkLightness{0.9}
\SetWatermarkScale{2}
\usepackage{hyperref}
\hypersetup{colorlinks=true,linkcolor=blue!50!black,urlcolor=blue!50!black}

\definecolor{revue}{RGB}{20,60,120}

\pagestyle{fancy}
\fancyhf{}
\renewcommand{\headrulewidth}{0.4pt}
\setlength{\headheight}{15pt}

% \entete{Matière}{Titre}{Type de document}
\newcommand{\entete}[3]{%
  \fancyhead[L]{\small\textsc{La Revue CAPES} -- #1}%
  \fancyhead[R]{\small #2}%
  \fancyfoot[C]{\thepage}%
  \thispagestyle{plain}%
  \begin{center}
    {\color{revue}\rule{\textwidth}{1.2pt}}\\[4pt]
    {\small\textsc{Concours directs CAP-CEG / CAPES -- Burkina Faso}}\\[6pt]
    {\LARGE\bfseries\color{revue} #2}\\[6pt]
    {\large\itshape #1 -- #3}\\[2pt]
    {\color{revue}\rule{\textwidth}{1.2pt}}
  \end{center}
  \vspace{0.5cm}%
}

% Encadré pour les remarques ajoutées dans les corrigés
\newcommand{\remarque}[1]{\par\medskip\noindent\fcolorbox{revue}{revue!6}{\parbox{\dimexpr\linewidth-2\fboxsep-2\fboxrule}{\textbf{\color{revue}Remarque.} #1}}\par\medskip}


\begin{document}
\entete{Mathématiques}{Sujet type n°1}{Énoncé}
\begin{spacing}{1.5}

\textbf{Exercice 1}

1. Soit $a$ et $b$ deux entiers naturels tels que $a \leq b$. On suppose que a divise $b$.

Montrer que a divise $b-a$. Quel est le quotient de $b-a$ par $a$?

2. En déduire les entiers naturels tels que $n-1$ divise $n+3$.

\textbf{Exercice 2 }

On définit l'application linéaire $f_a$ de $\mathbb{R}^3$ dans $\mathbb{R}^3$, $a$ étant un paramètre réel, par : 

\[
f_a : \begin{pmatrix}
x\\
y\\
z
\end{pmatrix} \longrightarrow \begin{pmatrix}
(2-a)x+(a-1)y\\
2(1-a)x+(2a-1)y\\
az
\end{pmatrix}
\]
 
1. Donner la matrice $M(a)$ représentant $f_a$. 

2. Trouver les noyaux K(0) de $f_0$ et K(1) de $f_1$. 

3. Dans la suite de l'énoncé, on prendra $a \neq 0$ et $a \neq 1$. Quel est le noyau K(a) de $f_a$ ? 

4. Montrer que $M(a)$ peut s'écrire sous la forme $A + aB$ où $A$ et $B$ sont deux matrices à déterminer. 

5. Calculer $AB, BA, A^2, B^2$; donner l'expression de $M(a)^2$ en fonction de $A, B$ et $a$. 

6. Calculer $[M(a)]^n$ en fonction de $a$, où $n$ est un entier naturel.

\textbf{Exercice 3 }

Soit un ensemble $E$ de $n$ éléments, $n\in \mathbb{N}$ ; on appelle combinaison d'ordre $p (0 \leq p \leq n)$ tout sous-ensemble de $E$ comportant $p$ éléments.  
Le nombre total de combinaisons d'ordre $p$ est égal à $C_{n}^{p}=\frac{n!}{p!(n-p)!}$.

1. Calculer $B = \sum_{p=0}^{p=n}C_{n}^{p}$
 
 2. Montrer que $C_{n}^{p}=f(n,p).C_{n-1}^{p-1}$ où $f(n,p)$ est une fonction de $n$ et $p$ à expliciter. 

3) Calculer $S = \sum_{p=0}^{p=n}p.C_{n}^{p}$.

\medskip\textbf{Exercice 4}

On considère l'équation différentielle
 
 $(E): y^{\prime\prime}+y^{\prime}+y=0$
 
 1. Résoudre l'équation différentielle $(E$) et étudier le comportement des solutions en $+\infty$ et $-\infty$

 2. Soit f une fonction de classe $C^2$ sur $\mathbb{R}$ telle que $f^{\prime\prime}+f^{\prime}+f$ soit $T$-périodique. Montrer que $f(x +T)-f(x)$ est solution de $(E)$. En déduire que si $f$ est bornée sur $\mathbb{R}, f$ est elle-même $T$-périodique.
 

\end{spacing}
\end{document}
